% Part: set-theory
% Chapter: z
% Section: story

\documentclass[../../../include/open-logic-section]{subfiles}

\begin{document}

\olfileid{sth}{z}{story}
\olsection{কাহিনিটির আরও বিস্তারিত বিবরণ}

\olref[story][approach]{sec}-এ সেট গঠনের প্রক্রিয়া সম্পর্কে
% BN-SRC-432: source spells Shoenfield's name as Schoenfield here.
শোয়েনফিল্ডের বিবরণ উদ্ধৃত করেছি। এখন কাহিনিটিকে একটু
নির্দিষ্ট করার জন্য আরও কয়েকটি নীতি লিখতে চাই। সেগুলি এই:
\begin{enumerate}
\item[] \stageshier. প্রত্যেক সেট কোনো-না-কোনো পর্যায়ে গঠিত হয়।
\item[] \stagesord. পর্যায়গুলি ক্রমবিন্যস্ত: কিছু পর্যায় অন্যগুলির
\emph{আগে} আসে।\footnote{আমরা আসলে—অপ্রকাশিতভাবে—ধরব যে
পর্যায়গুলি \emph{সুক্রমবিন্যস্ত}। এর অর্থ কী, তা
\olref[ordinals][]{chap}-এ ব্যাখ্যা করা হয়েছে। এটি যথেষ্ট
গুরুত্বপূর্ণ একটি অনুমান। বস্তুত, \citet{Scott1974}-এর একটি
অত্যন্ত চতুর কৌশলে এই অনুমান \emph{বাদ দিয়ে} পরে
\emph{নিষ্পন্ন} করা যায়। (এতে প্রথম পর্যায় আছে বলে
কেন ভাবব, তা-ও বোঝা যায়।) এখানে সে আলোচনা সম্ভব নয়;
আরও জানতে \citet{ButtonLT1} দেখো।}
\item[] \stagesacc. যেকোনো পর্যায়~$S$ এবং $S$-এর
\emph{আগে} গঠিত সেটগুলির যেকোনো সংগ্রহের জন্য, পর্যায়~$S$-এ
এমন একটি সেট গঠিত হয়, যার সদস্য ঠিক সেই সেটগুলি।
পর্যায়~$S$-এ অন্য কিছু গঠিত হয় না।
\end{enumerate}
এগুলি অনানুষ্ঠানিক নীতি। তবে এগুলি ব্যবহার করে আমরা
জার্মেলোর সেটতত্ত্বের কয়েকটি স্বতঃসিদ্ধের যথার্থতা দেখাতে
পারব।

(এখানে একটি সতর্কতা দরকার। আমরা পরে প্রচলিত নামযুক্ত
কয়েকটি সম্পূর্ণ প্রচলিত স্বতঃসিদ্ধ দেব; কিন্তু এইমাত্র
যে তির্যক হরফে লেখা নীতিগুলি দিয়েছি, সাহিত্যে সেগুলির
কোনো নির্দিষ্ট নাম নেই। আমরা কেবল সহায়ক হবে বলে
% BN-SRC-433: source omits the verb before “monikers.”
আশা করে কয়েকটি নাম দিয়েছি।)

\end{document}
